\documentclass[12pt]{article}
\usepackage[utf8]{inputenc}
\usepackage[T1]{fontenc}
\usepackage{hyperref}
\title{Mixtec, Tacahua Language Resources for K-12 Curriculum Development}
\author{Compiled by Claude Sonnet 5.5 (claude-sonnet-5-5) for Andr\'as Kornai}
\date{}
\begin{document}
\maketitle

\section{Introduction}
Tacahua Mixtec (ISO 639-3 xtt; Glottocode taca1257) is an Otomanguean language of Oaxaca, Mexico. It is also called Mixteco de Tacahua or Mixteco de Santa Cruz Tacahua. Wikipedia places it in the Atatl\'ahuca--San Miguel Mixtec group. The speaker figure of 28,000 on that page is for the whole group, not for Tacahua alone. The search found very little material specific to this variety, and no dictionary, grammar, primer or textbook.

\section{Primary Educational Resources}
No primers, textbooks or school programmes for Tacahua Mixtec turned up in the search.

\section{Linguistic and Reference Materials}
\begin{itemize}
\item Wikipedia, Atatl\'ahuca--San Miguel Mixtec (reached through the Tacahua Mixtec redirect). It gives classification information for the group: \url{https://www.wikipedia.com/wiki/Tacahua_Mixtec}
\end{itemize}

\section{Religious and Community Resources}
No Scripture translation or community organisation specific to xtt was found.

\section{Digital and Audio Resources}
\begin{itemize}
\item Global Recordings Network, Mixtec, Tacahua. It offers evangelistic and basic Bible-teaching audio recordings, with MP3 and slideshow downloads, aimed at people who are not literate: \url{https://www.globalrecordings.net/en/language/xtt}
\end{itemize}

\section{Curriculum Development Context}
Very little exists for this variety. A curriculum would have to be built from new documentation, ideally with the community, and could use the audio above as a source of spoken material. Related Mixtec varieties may have richer materials, but this search did not examine them.
\end{document}